\section{Methods}

\paragraph{Notation.} Time runs from noise ($t=0$) to data ($t=1$); $x_0\sim p_0=\mathcal N(0,I_d)$, $x_1\sim p_1$, $d=784$ for MNIST; a coupling $q(x_0,x_1)$ has these marginals; $v_\theta(x,t)$ is the network and $u_t(x)$ the marginal velocity field. Proofs are in Appendix~\ref{app:derivations}, further method details in Appendix~\ref{app:setup}.

\subsection{Conditional flow matching and its relation to diffusion}\label{sec:cfm}
CFM uses the straight conditional path $x_t=(1-t)x_0+tx_1$, whose velocity is $x_1-x_0$, and minimises
\begin{equation}\label{eq:cfm}
\mathcal L_{\mathrm{CFM}}(\theta)=\E_{t\sim U[0,1],\,(x_0,x_1)\sim q}\,\big\|v_\theta(x_t,t)-(x_1-x_0)\big\|^2 .
\end{equation}
\begin{proposition}[Appendix~\ref{app:d1}--\ref{app:d2}]\label{prop:cfm}
Assume $\E\|x_1\|^2<\infty$. Then $\mathcal L_{\mathrm{CFM}}(\theta)=\E_t\E\|v_\theta(x_t,t)-u_t(x_t)\|^2+C$ with $u_t(x)=\E[x_1-x_0\mid x_t=x]$ and $C=\E_t\E\operatorname{tr}\operatorname{Cov}(x_1-x_0\mid x_t)\ge0$ independent of $\theta$. Hence the minimiser is the posterior-mean velocity $u_t$, and if $u_t$ is locally Lipschitz with linear growth, $\dot x=u_t(x)$ started from $p_0$ reaches $p_1$ at $t=1$.
\end{proposition}
The constant $C$ depends on the coupling, so training losses under different couplings cannot be compared. I-CFM uses the independent coupling $p_0\otimes p_1$; OT-CFM solves an exact OT problem \citep{villani2009ot} inside each training batch with POT \citep{flamary2021pot}. With Gaussian noise and independent pairs, a diffusion path $x_t=a_tx_1+b_tx_0$ becomes the straight path after rescaling, $x_t/(a_t+b_t)$ at time $\tau=a_t/(a_t+b_t)$, and one deterministic DDIM step is exactly one Euler step of the I-CFM ODE on the grid $\tau_k$ (Proposition~\ref{prop:d9}). DDIM is therefore a fair reference; couplings that are not independent (OT, reflow) have no diffusion counterpart.

\subsection{Why few steps fail: one step lands on a posterior mean}\label{sec:onestep}
Since $x_1-x_0=(x_1-x_t)/(1-t)$, conditioning on $x_t$ gives
\begin{equation}\label{eq:postmean}
u_t(x)=\frac{\hat x_1(x,t)-x}{1-t},\qquad \hat x_1(x,t)=\E[x_1\mid x_t=x],
\qquad\text{hence}\qquad x+(1-t)\,u_t(x)=\hat x_1(x,t):
\end{equation}
one Euler step to $t=1$ lands exactly on the posterior mean (Proposition~\ref{prop:onestep}). From $t=0$ with independent pairs $\E[x_1\mid x_0]=\E[x_1]$, so every noise vector is sent to the data mean (Figure~\ref{fig:samples}). Few-step quality thus depends on how close the coupling is to a deterministic map $x_0\mapsto x_1$.

\subsection{Four ways to reduce the NFE}\label{sec:levers}
\paragraph{(i) Better pairs: minibatch OT, and its limit.} Relative to independent pairs a coupling saves $2\E\langle x_0,x_1\rangle$ of the pairing cost $\E\|x_0-x_1\|^2$, and the Gelbrich bound \citep{gelbrich1990} gives (Proposition~\ref{prop:bound})
\begin{equation}\label{eq:bound}
\text{saving}\;\le\;2\operatorname{tr}\Sigma^{1/2},\qquad \Sigma=\operatorname{Cov}(x_1),
\end{equation}
with equality for Gaussian data. On MNIST the independent cost is about 1,510 and $2\operatorname{tr}\Sigma^{1/2}=386$: no coupling can save more than 26\%, because most of $\E\|x_0\|^2=784$ is isotropic noise. We train with OT batches of 512 and 2,048 besides 128.
\paragraph{(ii) Where to train: the time distribution.} Sampling $t\sim\pi$ instead of $U[0,1]$ weights the regression error at time $t$ by $\pi(t)$ without changing the minimiser where $\pi>0$ (Proposition~\ref{prop:tdist}); with a finite network it decides where $v_\theta$ is accurate. We test logit-normal$(0,1)$ \citep{esser2024sd3} and $\pi=\tfrac12+\tfrac12\,\kappa/\!\int\kappa$, which trains more where the paths bend.
\paragraph{(iii) Straighter pairs: reflow.} Reflow \citep{liu2023rf} integrates the trained model from fresh noise $x_0$ to $\hat x_1$ and retrains on the pairs $(x_0,\hat x_1)$. For an exact field the marginals are preserved and the transport cost does not increase; the coupling is deterministic, so $\E[x_1\mid x_0]$ in \eqref{eq:postmean} is a single image (Proposition~\ref{prop:straight}). We use 60,000 dopri5 pairs per seed-0 model and fine-tune it for 50 epochs.
\paragraph{(iv) Better integration: a curvature-based time grid.} One Euler step of size $h$ has local error $\tfrac{h^2}{2}a(t)+O(h^3)$ with $a(t)=\tfrac{d}{dt}v(x(t),t)$ \citep[\S II.2]{hairer1993ode}. With $\kappa(t)=\E\|a(t)\|$ and a step density $n(t)$ with $\int_0^1n=N$, the summed local error is $E(n)\approx\int_0^1\kappa/(2n)\,dt$, and the Cauchy--Schwarz inequality gives (Proposition~\ref{prop:grid})
\begin{equation}\label{eq:sqrtk}
n^\star(t)=N\frac{\sqrt{\kappa(t)}}{\int_0^1\sqrt\kappa},\qquad
G=\frac{E(n^\star)}{E(\text{uniform})}=\frac{\big(\int_0^1\sqrt\kappa\big)^2}{\int_0^1\kappa}\le1 ,
\end{equation}
so $G$ predicts the gain before any sampling. We compare this $\sqrt\kappa$ grid with the uniform and shift grids \citep{esser2024sd3} for Euler, Heun and RK4 \citep{hairer1993ode}, with dopri5 \citep{dormand1980rk,chen2018node} as reference.

\subsection{Evaluation}\label{sec:eval}
\paragraph{FID and straightness.} FID \citep{heusel2017fid} uses Inception-v3 features \citep{szegedy2016inception} of 10,000 samples against the 60,000 MNIST \citep{lecun1998mnist} training images (clean-fid, \citealp{parmar2022cleanfid}); all samplers start from the same noise. NFE@20 is the NFE at which FID first reaches 20. Straightness is the path-length ratio $R=\E\big[\int_0^1\|\dot x\|\,dt\,/\,\|x(1)-x(0)\|\big]$, equal to 1 for straight paths.
\paragraph{Class coverage (the course method).} $k$-means \citep{lloyd1982kmeans} with $k=10$ is fitted on 30,000 real images. The cluster shares $\hat p_k$ of $n$ generated images are compared with the shares $p_k$ of $m=30{,}000$ held-out real images by
\begin{equation}\label{eq:tv}
\mathrm{TV}=\tfrac12\sum_k|\hat p_k-p_k|,\qquad
\E[\mathrm{TV}]\big|_{\text{perfect model}}\approx\tfrac12\sum_k\sqrt{\tfrac{2}{\pi}\,p_k(1-p_k)\big(\tfrac1n+\tfrac1m\big)}=0.014,
\end{equation}
the second expression being the noise floor (Proposition~\ref{prop:tvfloor}), and by Pearson's two-sample $\chi^2$ test of homogeneity \citep{agresti2013cda}, needed because $p_k$ is itself estimated (Appendix~\ref{app:course}).
\paragraph{Decision rule.} With FID $F_{s,r}=\mu+a_s+e_{s,r}$ for training seed $s$ and noise draw $r$, separately trained models must differ by more than $2\hat\sigma_{\text{seed}}$ (three seeds). Samplers on one model are paired, so $a_s$ cancels; their differences must exceed the sampling noise and keep their sign on the other seeds. With three seeds we report sign agreement, not significance.
